% xelatex
\documentclass{resume}
\usepackage{setspace}
\usepackage{hyperref}
\usepackage{titlesec}
\usepackage{ctex}
\usepackage{xeCJK}
% \setCJKmainfont{Microsoft YaHei}

\titleformat*{\subsection}{\normalsize\bfseries}
\hypersetup{
    colorlinks=true,
    linkcolor=black,
    filecolor=magenta,
    urlcolor=blue,
}

\begin{document}
\linespread{1.2}
\pagenumbering{gobble} % suppress displaying page number


\name{段浩东}

\basicInfo{
  \email{dhd.efz@gmail.com} \textperiodcentered
  \faHome\ \href{http://kennymckormick.github.io}{主页}
}

\vspace{-2mm}

\section{教育经历}

\datedsubsection{\textbf{香港中文大学}}{2019 -- 至今}
研究视频理解相关课题，导师为\href{https://dahua.me}{林达华教授} \\
信息工程系，博士生

\datedsubsection{北京大学}{2015 -- 2019}
GPA 3.77/4.00, rank 1st in Data Science students \\
元培学院，数据科学方向，本科生

\vspace{-2mm}

\section{研究方向}
博士期间主要研究计算机视觉中的视频理解任务，特别对基于人体骨架的高效（数据高效 \& 计算高效）视频理解感兴趣。

\vspace{-2mm}

\section{科研项目}
\datedsubsection{Unified Framework for Skeleton Action Recognition in the Wild}{2022}
Design a unified framework \textbf{SkeleTR} to handle instance-level, group-level, video-level skeleton action recognition in the wild, which outperforms the GCN SOTAs by a large margin on all benchmarked tasks.

\vspace{-2mm}

\datedsubsection{Skeleton Action Recognition with Dynamic Group-wise GCN, \href{https://arxiv.org/abs/2210.05895}{Paper}, \href{https://github.com/kennymckormick/pyskl/tree/master/configs/dgstgcn}{Code}}{2021 -- 2022}
Design a dynamic GCN for skeleton action recognition,
that requires no pre-defined skeleton topology.
It achieves strong recognition performance, surpasses previous SOTAs on NTURGB+D and K400.

\vspace{-2mm}

\datedsubsection{Video Self-supervised Learning via Ranking-based Transformation Recognition, \href{https://arxiv.org/abs/2205.02028}{Paper}}{2021}
Show the great potential of RecogTrans (recognizing transformations applied to video clips) video SSL by introducing a unified Ranking-based formulation.
The proposed method significantly outperforms previous RecogTrans approaches on action recognition (UCF Top1 +6\%) and video retrieval (UCF R@1 +20\%).

\vspace{-2mm}

\datedsubsection{Efficient Video Recognition for Untrimmed Videos, \href{https://arxiv.org/abs/2201.04388}{Paper}, \href{https://github.com/MCG-NJU/OCSampler}{Code}}{2021}
Propose OCSampler for untrimmed video recognition.
It samples frames from candidates to form one representative clip.
The framework (w. R50) achieves 82.2\% Top-1 with 1-clip testing (only 52 GFLOPs / video).

\vspace{-2mm}

\datedsubsection{Skeleton Action Recognition with 3D ConvNets, \href{https://arxiv.org/abs/2104.13586}{Paper}, \href{https://github.com/kennymckormick/pyskl/blob/main/tools/data/data_doc.md}{Dataset}, \href{https://github.com/kennymckormick/pyskl/tree/master/configs/posec3d}{Code}}{2020 -- 2021}
Devise a novel 3D-ConvNet based paradigm (\textbf{PoseC3D}: 2D keypoint heatmaps $\rightarrow$ 3D heatmap volumes $\rightarrow$ 3D-CNN recognizer) for skeleton action recognition. PoseC3D outperforms previous skeleton action recognition approaches by a considerable margin across various benchmarks (NTURGB+D, Kinetics, \emph{etc.}).

\vspace{-2mm}

\datedsubsection{Mitigating Unwanted Bias in Action Recognition, \href{https://arxiv.org/pdf/2209.09393.pdf}{Paper}, \href{https://github.com/kennymckormick/ARAS-Dataset}{Dataset}}{2020}
Demonstrate that deep learning based video recognition is biased towards factors like scene / objects.
Create a quantitative benchmark to evaluate the bias, mitigate the bias with adversarial training and diversified web data.

\vspace{-2mm}

\datedsubsection{Omni-sourced Webly-supervised Video Recognition, \href{https://arxiv.org/abs/2003.13042}{Paper}, \href{https://github.com/open-mmlab/mmaction2/tree/master/tools/data/omnisource}{Dataset}, \href{https://github.com/open-mmlab/mmaction2/tree/master/configs/recognition/omnisource}{Code}}{2019 -- 2020}
Propose \textbf{OmniSource} for webly supervised video recognition, which utilizes various kinds of web data, including images, trimmed / untrimmed videos for trimmed video recognition. Achieve 83.6\% Top-1 on Kinetics400.

\vspace{-2mm}

\datedsubsection{Triplet Representation for Human Body,  \href{http://openaccess.thecvf.com/content_ICCV_2019/papers/Duan_TRB_A_Novel_Triplet_Representation_for_Understanding_2D_Human_Body_ICCV_2019_paper.pdf}{Paper},  \href{https://github.com/kennymckormick/Triplet-Representation-of-human-Body}{Dataset}} {2018 -- 2019}
Design a triplet representation named TRB (and its estimation method) to represent 2D human body, which includes both human pose and shape information. The representation can be used in human shape editing.

\vspace{-2mm}

\section{实习经历}

\datedsubsection{实习应用科学家, AWS AI}{2022}
研究一般场景下的骨骼动作识别, 由 \href{https://xumingze0308.github.io/}{徐铭泽} 博士和 \href{https://alessandrobergamo.com/}{Alessandro Bergamo} 博士指导。

\vspace{-2mm}

\datedsubsection{实习研究员, 上海人工智能实验室}{2021 -- 2022}
开发维护开源算法库 \href{https://github.com/open-mmlab/mmaction2}{MMAction2}, 由 \href{https://chenkai.site/}{陈恺} 博士指导。

\vspace{-2mm}

\datedsubsection{实习研究员, 商汤科技}{2017 -- 2019}
研究人体姿态估计 (骨骼及轮廓关键点), 由 \href{https://scholar.google.com/citations?user=KZn9NWEAAAAJ}{刘文韬} 博士指导。

\vspace{-2mm}



\section{预印本}

\textbf{Haodong Duan}, Jiaqi Wang, Kai Chen, Dahua Lin \\
\emph{DG-STGCN: Dynamic Spatial-Temporal Modeling for Skeleton-based Action Recognition} \\

\vspace{-2mm}

\section{发表论文}

Yujie Zhou, \textbf{Haodong Duan}, Anyi Rao, Bing Su, Jiaqi Wang \\
\emph{Self-supervised Action Representation Learning from Partial Spatio-Temporal Skeleton Sequences} (AAAI 2023) \\

\textbf{Haodong Duan}, Yue Zhao, Kai Chen, Yuanjun Xiong, Dahua Lin\\
\emph{Mitigating Representation Bias in Action Recognition: Algorithms and Benchmarks} (ECCVW 2022) \\

\textbf{Haodong Duan}, Jiaqi Wang, Kai Chen, Dahua Lin\\
\emph{PYSKL: Towards Good Practices for Skeleton Action Recognition} (MM 2022) \\

\textbf{Haodong Duan}, Yue Zhao, Kai Chen, Dahua Lin, Bo Dai\\
\emph{Revisiting Skeleton-based Action Recognition} (CVPR 2022 Oral) \\

\textbf{Haodong Duan}, NanXuan Zhao, Kai Chen, Dahua Lin\\
\emph{TransRank: Self-supervised Video Representation Learning via Ranking-based Transformation Recognition} (CVPR 2022 Oral) \\

Jintao Lin, \textbf{Haodong Duan}, Kai Chen, Dahua Lin, Limin Wang \\
\emph{OCSampler: Compressing Videos to One Clip with Single-step Sampling} (CVPR 2022) \\

\textbf{Haodong Duan},  Yue Zhao, Yuanjun Xiong, Wentao Liu, Dahua Lin \\
\emph{Omni-sourced Webly-supervised Learning for Video Recognition} (ECCV 2020)  \\

\textbf{Haodong Duan}, Kwanyee Lin, Sheng Jin, Wentao Liu, Chen Qian,  Wanli Ouyang\\
\emph{TRB: A Novel Triplet Representation for Understanding 2D Human Body} (ICCV 2019)  \\

\vspace{-2mm}

\section{开源项目}
为开源项目 \href{https://github.com/open-mmlab/mmaction}{MMAction},  \href{https://github.com/open-mmlab/mmaction2}{MMAction2}, \href{https://github.com/kennymckormick/pyskl}{PYSKL} 的主要开发者和维护者。

\vspace{-2mm}

\section{研究社群服务}
会议审稿: ICCV[21-23], AAAI[22-23], CVPR[22-23], ECCV22, NeurIPS22, WACV23, ICML23.  \\
期刊审稿: TCSVT, SPL, JVCIR.

\vspace{-2mm}

\section{语言技能}
\begin{itemize}[parsep=0.5ex]
    \item TOEFL iBT: 104pt (Reading: 30, Listening: 28, Speaking: 20, Writing: 26)
    \item GRE: 322pt (Verbal: 152, Quantitative: 170)
\end{itemize}

\end{document}
